% Generated draft: requires mathematical and source review.
\documentclass[11pt]{article}
\usepackage[margin=25mm]{geometry}
\usepackage{fontspec}
\setmainfont{Latin Modern Roman}
\usepackage{amsmath,amssymb,mathtools,unicode-math}
\setmathfont{Latin Modern Math}
\usepackage{enumitem,needspace,xurl,hyperref}
\setlength{\parindent}{0pt}
\setlength{\parskip}{5pt}
\setlength{\emergencystretch}{3em}
\newcommand{\N}{\mathbb{N}}
\newcommand{\Z}{\mathbb{Z}}
\newcommand{\Q}{\mathbb{Q}}
\newcommand{\R}{\mathbb{R}}
\newcommand{\C}{\mathbb{C}}
\newcommand{\F}{\mathbb{F}}
\newcommand{\A}{\mathbb{A}}
\newcommand{\PP}{\mathbb{P}}
\newcommand{\E}{\mathbb{E}}
\newcommand{\eps}{\varepsilon}
\newcommand{\dd}{\,\mathrm{d}}
\begin{document}
\section*{UnsolvedMath ID: 11000244}
\textbf{Source problem code:} catalogue entry
\subsection{Definitions and mathematical statement}
Let \(M=\mathbb H^3/\Gamma\) be a complete finite-volume hyperbolic three-manifold. The principal question is whether \(\Gamma\) contains a closed surface subgroup:
\[\exists h\geq2,\qquad\pi_1(\Sigma_h)\hookrightarrow\Gamma,\]
where \(\Sigma_h\) is a closed connected orientable genus-\(h\) surface. This requires an injective homomorphism, not an embedded surface in \(M\), nor merely the free fundamental group of a surface with boundary. Orientation may be handled by the double cover as the source notes.

The historical source says [8] answers the noncompact finite-volume case, leaving the closed case then far from resolution. The supplied later summary reports the closed case solved by Kahn--Markovic (2012); combined with the earlier noncompact result this is its reported resolution of the stated finite-volume question. The historical difficulty statement is not a present-status assertion.

The following discussion reports that a closed essential surface in a noncompact finite-volume \(M\) is either quasi-Fuchsian or has accidental parabolics, and that these surfaces are geometrically finite in this setting. Quasi-Fuchsian refers to the surface-group representation; its nontrivial elements are loxodromic (called hyperbolic in the source). An accidental parabolic is a nontrivial surface-group element sent to a parabolic ambient isometry. The source says [8] constructs the latter type only, and poses existence of closed quasi-Fuchsian surfaces as a further historical open question, mentioning progress by Masters--Zhang [27]. That separate qualification is not automatically settled by the summary's closed-manifold result.

Its description of [8] is retained: pass to an orientable cover if needed and use residual finiteness to obtain a finite cover \(\widehat M\) whose compact core has at least three torus boundary components. The source reports \(b_1(\widehat M)\geq3\), and an embedded orientable essential surface \(F\) with nonempty boundary missing at least one such component. In a further finite cyclic cover, two copies of \(F\) can be tubed along their boundaries to form a closed embedded essential surface of genus at least two. Its projection to \(M\) yields the desired subgroup, possibly as an immersed surface. The tubed surface contains accidental parabolics corresponding to peripheral elements from \(\partial F\). These are supplied construction reports; exact hypotheses and the proof are in the absent reference.

The final subsection describes combination-method context. Suppose \(G_0,G_1,G_2\leq\Gamma\), \(G_1\cap G_2=G_0\), and \(\Gamma\) acts on a set \(X\). A proper interactive pair \((\Theta_1,\Theta_2)\) satisfies:
\begin{enumerate}
\item \(\Theta_i\ne\varnothing\) for \(i=1,2\).
\item \(\Theta_1\cap\Theta_2=\varnothing\).
\item \(G_0\) preserves each \(\Theta_i\).
\item For every \(g_1\in G_1\setminus G_0\), \(g_1\Theta_2\subseteq\Theta_1\); for every \(g_2\in G_2\setminus G_0\), \(g_2\Theta_1\subseteq\Theta_2\).
\item For each \(i=1,2\) there exists \(\theta_i\in\Theta_i\) such that \(\theta_i\notin g_i\Theta_{3-i}\) for every \(g_i\in G_i\setminus G_0\).
\end{enumerate}
The source says repeated Maskit combination arguments provide another viewpoint on [8], referring to [23] and [26], Chapter VII. It announces a following combination theorem but does not supply it. No unstated amalgam or discreteness theorem is inserted, and the undefined contextual \(\Gamma_g\) is not identified with a new object.
\subsection{Short English statement}
Every finite-volume hyperbolic three-manifold group should contain a closed genus-at-least-two surface group; the record reports this resolved using the earlier cusped result and Kahn--Markovic's closed case. Its accidental-parabolic construction and proper-interactive-pair context are preserved separately.
\subsection{Sources}
\begin{itemize}
\item Local UnsolvedMath registry, permanent ID \(11000244\), exact source code \texttt{catalogue entry}.
\item Farb, \emph{Problems on Mapping Class Groups and Related Topics}, Question 2.1, PDF page 282, with subsequent Sections 2.2--2.3; the supplied running footer names A. W. Reid: \url{https://www.math.uchicago.edu/~farb/papers/mcgbook.pdf}. Registry access 2026-07-29; pdf-font-structure extraction.
\item Internal references [8] (noncompact construction), [7] (related discussion), Masters--Zhang [27], and Maskit references [23], [26], Chapter VII. Full bibliographic entries are absent.
\item The summary attributes the closed-case solution to Kahn--Markovic (2012, \emph{Annals of Mathematics}), mentioning ``The good pants homology and the Ehrenpreis conjecture'' and ``Immersing almost geodesic surfaces...'' together. These distinct/truncated titles and publication details require precise bibliographic matching; no missing title is completed here.
\end{itemize}
\end{document}
